\documentclass[10pt,a4paper]{amsart}
\usepackage[margin=19mm]{geometry}
\usepackage{amsmath,amssymb,mathtools}
\usepackage[scr=rsfs,cal=euler]{mathalfa}
\usepackage{xfrac}
\usepackage[unicode,hidelinks,bookmarks=false]{hyperref}
\newcommand{\ie}{i.e.\ }
\newcommand{\cf}{cf.\ }
\newcommand{\resp}{respectively\ }
\newcommand{\Subsection}{\subsection}
\providecommand{\nobreakdash}{\nobreak\hbox{-}}
\setlength{\emergencystretch}{2em}
\multlinegap0pt
\usepackage{bm}
\usepackage{bbm}
\usepackage{dsfont}
\usepackage{xspace}
\usepackage{cancel}
\usepackage{mathtools}
\usepackage{amsthm}
\makeatletter
\@ifundefined{thm}{\newtheorem{thm}{Thm}}{}
\@ifundefined{lemma}{\newtheorem{lemma}[thm]{Lemma}}{}
\makeatother
\def\cA{{\mathcal A}}
\def\mbbR{\mathbb{R}}
\def\mbbS{\mathbb{S}}
\def\mrH{\textrm H}
\def\mrd{\textrm d}
\def\rmd{\mathrm d}
\newcommand{\beq}{\begin{equation}}
\newcommand{\eeq}{\end{equation}}
\title[Gradient Flows of Interfacial Energies: Curvature Agents and]{Gradient Flows of Interfacial Energies: Curvature Agents and Incompressibility}
\author{Keith Promislow, Truong Vu, Brian Wetton}
\date{Source: arXiv:2509.07380v1 (CC BY 4.0)}
\begin{document}
\maketitle

\noindent\emph{Source and licence.} Gradient Flows of Interfacial Energies: Curvature Agents and Incompressibility. Version arXiv:2509.07380v1 (CC BY 4.0), distributed under the Creative Commons Attribution 4.0 International licence, as recorded in the arXiv licence metadata for this exact version and shown on the abstract page https://arxiv.org/abs/2509.07380v1. Full rights notice: licenses/arxiv-2509.07380-CC-BY-4.0.txt.\par\smallskip

\noindent\textit{Editorial scope.} Counted unit: the Lemma whose source label is \texttt{l:H\_coerc} in the article (source0.tex lines 1346--1354, source label \texttt{l:H\_coerc}), delivered together with the article's own complete original proof of it (source0.tex lines 1355--1414, one \texttt{proof} environment of 60 lines). No mathematics is written, repaired or re-derived: statement and proof are the article's own text line for line. Required context retained uncounted: none. Every definition, display and label of those lines is retained unchanged. Omitted from this package and not redistributed: 1--1345, 1415--1575. Nothing inside a retained excerpt is omitted or altered: every retained body is delivered exactly as the article's lines are written, so no omission receipt of any kind accompanies this package. Citations: the retained excerpts carry no citation command at all, so no bibliography entry of the article is transcribed and no reference is redistributed. Reference adaptation: none was needed and none was made; every reference inside the delivered unit resolves inside it, so no unresolved reference or citation is produced. Printed slips kept literal: no printed word, symbol, hypothesis or number of the counted statement or of its proof was altered. Layout and class: the article's own class is \texttt{amsart} with packages including amssymb, amsfonts, stmaryrd, fancyhdr, amsmath, amsthm, amsbsy, titletoc, latexsym, amscd, wrapfig, multirow, amsaddr, comment; the wrapper is the lane's pinned \texttt{amsart} editorial wrapper at 10pt on A4, which restates the theorem-like environments the retained text needs and adds the article's own macro definitions that the retained text needs (\texttt{\textbackslash beq}, \texttt{\textbackslash cA}, \texttt{\textbackslash eeq}, \texttt{\textbackslash mbbR}, \texttt{\textbackslash mbbS}, \texttt{\textbackslash mrH}, \texttt{\textbackslash mrd}, \texttt{\textbackslash rmd}), with extra wrapper packages (bm, bbm, dsfont, xspace, cancel, mathtools). The delivered numbering is the unit's own. Assets: no retained span contains a figure environment, a graphics-inclusion command, a PDF-page inclusion, a table or a TikZ picture; no image file is copied into this package, the package declares no asset dependency and no image file is redistributed with it. Licence: version arXiv:2509.07380v1 of this article is distributed under the Creative Commons Attribution 4.0 International licence, as recorded in the arXiv licence metadata for this exact version and shown on the abstract page https://arxiv.org/abs/2509.07380v1. A full rights notice is delivered at licenses/arxiv-2509.07380-CC-BY-4.0.txt. Characters: every retained excerpt is pure ASCII, so the delivered engine, XeLaTeX with its default Latin Modern text font, reports no missing character.
\begin{lemma}
\label{l:H_coerc}
For each $U\in\cA$ the spectrum of $\mrH(U)$ is real and non-negative. Fix $M>0$ in $\mbbR$ and form  the set $\cA_M\subset\cA$ comprised of $U$  whose image $\Gamma_U$ satisfies the length bound $|\Gamma_U|=\int_\mbbS \mrd\sigma \leq M.$  Then the spectrum of $\mrH(U)$ is uniformly bounded away from zero, that is there exists $\nu_1=\nu_1(M)$ given in \eqref{e:nu0-def} such that
\beq\label{e:H-L2coer}
\langle \mrH u,u\rangle \geq \nu_1 \|u\|_{H^1}^2,
\eeq
for all $U\in\cA_M$ and all $u\in H^1(\mbbS).$
In particular the spectrum satisfies  for all $U\in\cA_M.$ 
\end{lemma}

\begin{proof}
Since $\mrH$ is self adjoint and its bilinear form is nonnegative its spectrum is also real and non-negative. 
To establish a uniform lower  bound on its spectrum over $\cA_+$ we use the constraint $\int_\mbbS \kappa\mrd \sigma= 2\pi$. 
From Young's inequality we have
$$\begin{aligned}
2\pi =\int_\mbbS \kappa  \rmd \sigma &\leq \left(\int_\mbbS\kappa^2\mrd\sigma\right)^\frac12 \left(\int_\mbbS\mrd\sigma\right)^\frac12,\\
\end{aligned}
$$
from which we deduce the lower bound
$$
\| \kappa \|_{L^2} \geq  \frac{2\pi}{\sqrt{M}}.
$$
The bilinear form induced by $\mrH$ takes the form
$$\begin{aligned}
    \langle \mrH u, u\rangle_\mbbS&=\int_\mbbS |\nabla_s u|^2+\kappa^2u^2 \,\mrd \sigma.
 %   &\geq  \frac{g_-}{g_+^2}\int_\mbbS |\partial_s u|^2+g_+^2\kappa^2 u^2\mrd s.
\end{aligned}$$
 Define the set 
 $$S_\kappa:=\left\{s\in\mbbS\,\bigl|\,\kappa^2>\frac{2\pi^2}{M^2}\right\}.$$ 
 Since $|\mbbS|=1$ the lower bound on the $L^2$ norm of $\kappa$ implies that
\beq\label{e:Sk-lower}\int_{S_\kappa}\kappa^2\,\mrd\sigma \ge \frac{4\pi^2}{M}-\frac{2\pi^2}{M^2}\int_{\mbbS\backslash S_\kappa} \mrd\sigma \geq \frac{2\pi^2}{M}.\eeq
Consider a function $u\in H^1(\mbbS)$
with $\|u\|_{L^2}=1$ and
$$  \langle \mrH u, u\rangle_\mbbS\leq \nu,$$
for some $\nu>0.$
This implies that
$$ \int_{S_\kappa} u^2\mrd \sigma<\frac{\nu M^2}{2\pi^2}.$$
Let $u_\kappa\in\mbbR$ denote the infimum of $|u|$ over $S_\kappa.$ Using \eqref{e:Sk-lower} we have
$$ \nu\geq \int_{S_\kappa} |\nabla_s u|^2+\kappa^2 u^2\mrd \sigma \geq u_\kappa^2 \int_{S_\kappa}\kappa^2\mrd\sigma\geq  \frac{2\pi^2}{M} u_\kappa^2,$$
and hence
$$ u_\kappa\leq \sqrt{\frac{\nu M}{2\pi^2}}.$$
The set $S_u:=\{s\in\mbbS\,\bigl|\, u^2>\frac{1}{2M}\}$, is not empty. Indeed, since $\|u\|_{L^2}=1$ we have the lower bound on its size 
$$ \int_{S_u} u^2\mrd \sigma=1-\int_{\mbbS\backslash S_u}u^2\rmd \sigma \geq \frac12.$$
Fix an  arc-length function $g$. If $u$ takes values $u_\pm\in\mbbR$ with $u_+>u_-$ over a separation of $\ell$, then a standard calculus of variations argument on the quantity
$$ \int_0^\ell |\nabla_s u|^2\mrd \sigma,$$
for $u\in H^1([0,\ell])$ subject to $u(0)=u_-$ and $u(\ell)=u_+,$ shows that the minimum is achieved at
$$ u_*(s) =\frac{u_+-u_-}{\int_0^\ell\mrd\sigma}\int_0^s\mrd\sigma.$$
This minimizer has constant surface gradient
$$\nabla_s u_* =\frac{u_+-u_-}{\int_0^\ell\mrd\sigma}.$$
Consequently if any $u\in H^1(\mbbS)$ attains the values $u_\pm$ over a separation $\ell,$ we have the lower bound
\beq\label{e:grad-bnd}
\int_\mbbS |\nabla_s u|^2\rmd \sigma >\frac{(u_+-u_-)^2}{\int_0^\ell \mrd\sigma}\geq \frac{(u_+-u_-)^2}{M}.
\eeq
If $\nu< \pi^2/M$ then we have  $u_\kappa<\frac{1}{\sqrt{2M}}$ and may takes these values as our $u_\pm$. This affords a lower bound on the gradient integral and hence on $\nu$. Specifically \eqref{e:grad-bnd} yields
$$ \nu \geq \langle \mrH u, u\rangle > \frac{1}{M} \left(\frac{1}{\sqrt{2M}}-\sqrt{\frac{\nu M}{2\pi^2}}\right)^2> \frac{1}{2M^2}\left(1-\frac{M}{\pi}\sqrt{\nu}\right)^2.$$
This inequality can only hold if $\nu$ satisfies a universal lower bound 
\beq\label{e:nu0-def}\nu\geq\nu_0:= \frac{c}{M^2},
%\left(\frac{\pi}{\pi\sqrt{2}+1}\right)^2\frac{1}{M^2}.
\eeq
for a constant $c.$ This  establishes the uniform $L^2$ coercivity of $\mrH$ over $\cA_M$. To obtain $H^1$ coercivity we observe
that 
$$\begin{aligned}
    \langle \mrH u, u\rangle &\geq \frac 12 \int_\mbbS |\nabla_s u|^2\mrd\sigma+ \frac12\langle\mrH u, u\rangle,\\ 
    &\geq \frac12\int_\mbbS |\nabla_s u|^2\mrd\sigma+   \frac {\nu_0}{2} \|u\|_{L^2}^2.
    \end{aligned} 
 $$
 Taking $\nu_1=\frac12 \min\{1,\nu_0\}$ 
% $$\nu_1=\min\left\{\frac12, \frac12 \left(\frac{\pi}{\pi\sqrt{2}+1}\right)^2\frac{1}{M^2} \right\},$$
 yields the $H^1$ coercivity.
\end{proof}

\end{document}
